
This paper presented a design of a modular holonomic hexarotor robot
for microgravity environments. The propulsion is based on 6 propellers
oriented in such a way that both holonomy is attained and the minimum
upper bound of the thrust across all directions is maximized. 
% The design is modular, so that maintenance is simplified. We preview as potential applications telepresence and object transportation. The robot can also be used to host microgravity experimental test-beds.

We used a multi-criteria optimization approach to guide the design parameters of the proposed propulsion system. We also include a convergent controller for the vehicle, and a validation of our design in a realistic physics-based simulator. Simulation results show the vehicle navigating between waypoints. In particular, we show that the convergence to the waypoints is robust to localization noise and to unmodeled dynamics, such as attachment to heavy loads.

A real implementation of Space CoBot will face several challenges. One
of them is the asymmetry of propeller trust depending on the rotation
direction can be addressed by appropriately scaling the $u$ to
rotation speed mapping. Other issues include modeling the motor and
propeller inertia and the quadratic relation between propeller speed
and thrust\footnote{The lower the propeller speed is, a higher change
  in its speed is needed to achieve the same change in force.}.

Future work will address: (1)~vision-based navigation using the onboard camera array, (2)~mobile dexterous manipulation with the onboard robotic arm, encompassing both single robot and cooperative multirobot, and (3)~robustify the controller to better handle unmodeled dynamics, such as heavy loads. We will also aim at constructing a prototype for validating the design both in 2D motion, e.g., frictionless table on Earth, and in microgravity testbeds, e.g., in a parabolic flight aircraft.

%%% Local Variables:
%%% mode: latex
%%% TeX-master: "main"
%%% End:
